% ============================================================
% front/abstract.tex
% ============================================================
\chapter*{Abstract}
\addcontentsline{toc}{chapter}{Abstract}

AG Technologies, a Cameroonian company specializing in digital
solutions, operated a fleet of physical servers hosting critical
applications, managed manually and separately, with no single point of
control and no consistent traceability of actions performed. This
observation led to the design of AGT Infra, an internal platform that
centralizes the management of this fleet from a single entry point.

AGT Infra positions itself as an orchestrator rather than an engine: it
drives already proven building blocks (Docker, SSH, GitHub, Nginx,
Certbot) without ever reimplementing them, relying on a dynamic access
control model (RBAC/IBAC) in which an explicit denial always takes
precedence over a broader grant, together with systematic traceability
of every sensitive action. The tool covers server and Docker container
management, continuous deployment from Git repositories, secure
application publication under HTTPS, autonomous monitoring through a
dedicated agent, and an integrated web console.

Most of these domains are now implemented and verified under real
conditions. Automated backups, automatic scaling, and an active security
engine remain identified needs that have not yet been implemented, and
constitute the main areas for future development of the project.

\bigskip
\noindent\textbf{Keywords:} infrastructure management, DevOps, dynamic
access control (RBAC/IBAC), continuous deployment, containerization,
monitoring, SSH orchestration.